\documentclass[11pt,a4paper]{report}

\usepackage{microtype}
\usepackage{xcolor}
\usepackage{subcaption}
\usepackage{amsmath}

% Editing macros -- delete for the final draft (matches project convention in
% 1_Chapter_Intro/a_11th_aug_introdraft.tex and 5_Chapter_results_evaluation/b_draft_plan_14th_aug.tex)
\newcommand{\fc}[1]{\textcolor{red}{[FactC: #1]}}
\newcommand{\ct}[1]{\textcolor{orange}{[CITE: #1]}}
\newcommand{\app}[1]{\textcolor{blue}{[APP: #1]}}
\newcommand{\sr}[1]{\textcolor{teal}{[Sreword: #1]}}
\newcommand{\fig}[1]{\textcolor{purple}{[FIG: #1]}}
\newcommand{\ai}[1]{\textcolor{brown}{[AI: #1]}}
\newcommand{\chk}[1]{\textcolor{orange}{[CHECK: #1]}}

\begin{document}

% ============================================================================
% WORKING NOTES -- DELETE THIS WHOLE BLOCK BEFORE FINAL SUBMISSION.
% Everything from here to "END WORKING NOTES" is planning scaffolding, not
% dissertation prose. The actual chapter draft starts at \chapter{Results and
% Discussion} below. Source material: c_draft_plan_15th_1149pm_audit.md (all
% numbers), d_draft_plan_16th_855pm_plotting.md (figure catalogue),
% 1_Chapter_Intro/a_11th_aug_introdraft.tex (research objectives),
% 4_Chapter_methodology/h_methodology_2026_08_15_1117_my_own_edits_to_g.tex
% (model/evaluation definitions). No new numbers were computed for this file;
% every value below is copied from c_'s already-audited tables.
% ============================================================================

\section*{1. Argument outline per research question}

\textbf{RQ1 (top-height prediction, Objective 3).} A tuned XGBoost model beats every
neural architecture on both cohorts, and the margin holds fold-by-fold, not just on
average. Chapman-Richards guidance does not close this gap at any tested loss weight;
removing the guidance recovers accuracy but destroys a specific interpretability
property (asymptote/growth-rate separability) that only the guided architecture has to
begin with. The spatial-block split matters: an easy plot-level split inflates the
unconstrained DNN's score by up to 0.197 R2 while barely moving the guided models,
which is the direct evidence for treating spatial-block as the primary split
throughout the rest of the chapter.

\textbf{RQ2 (attribution to the shared curve, Objective 4).} A flexible model
conditioned on environmental data is not a uniform improvement: it corrects the plots
the shared curve already fits worst and is noticeably worse on the plots it already
fits well. A permutation control shows only the first half is a genuine effect of
real environmental data; the second half appears at a similar size even with fake,
scrambled environmental values, so it is a cost of model flexibility in general, not
of environmental conditioning specifically. A plain XGBoost model also reproduces the
overall pattern from the same real features, which points to the environmental data
itself, not to any physics-guided architecture, as the source of the genuine (Q4)
correction. Once the persistent (whole-history) departure is isolated and modelled
directly, canopy cover and thinning history dominate every method tested, slope is the
only true abiotic variable that agrees in sign and size across model families, and a
large share of the departure's spatial clustering survives every environmental set
tested, which motivates RQ3's spatially-varying approach directly.

\textbf{RQ3 (plot-specific curve deviation, Objective 5).} A spatially-varying model
(GNNWR) shows a consistent point-estimate edge over global attribution models on the
cohort with enough spatial density to support it, corroborated by an independent
reduction in residual spatial clustering, but this edge is not confirmed by the
strictest confidence-interval bar and the effect does not transfer to the smaller
cohort. Separately, the plots with the most extreme fitted deviations trace, for about
half of them, to how the target itself is constructed under a fixed growth-rate
assumption rather than to a genuine ecological anomaly, a mechanism confirmed by
tracing the closed-form fit directly rather than inferred from the pattern alone.

\section*{2. Ranked headline findings}

\begin{enumerate}
\item RQ1: tuned XGBoost outperforms every neural model on both cohorts, fold-by-fold,
  not only on average.
\item RQ2b: the shared-curve residual keeps substantial spatial clustering
  (Moran's I 0.65--0.74) regardless of how much environmental data is added.
\item RQ2a: a flexible model trades accuracy between subgroups regardless of what it
  is fed, but only the correction on badly-fit plots is a genuine effect of real
  environmental data; the accuracy loss on well-fit plots is not, confirmed by a
  permutation control.
\item RQ3: GNNWR's accuracy edge on 4survey is corroborated by an independent
  reduction in residual clustering, but is not CI-clean and does not hold on 6survey.
\item RQ2b: canopy cover and thinning history dominate attribution across NLME,
  Elastic Net and XGBoost alike.
\item RQ3: roughly half of the most extreme fitted-asymptote outliers trace to the
  target's own fixed-growth-rate construction, not to disturbance.
\item RQ1: the physics constraint trades real accuracy for a specific, narrowly-scoped
  parameter-identifiability property in the rate-conditioned guided model.
\item RQ2a: a plain XGBoost model reproduces RQ2a's residual-reduction pattern almost
  exactly, pointing to the data, not the architecture, as the cause.
\item RQ3: on 6survey, XGBoost's own attribution disagrees with EN and GNNWR on which
  variable ranks first, and XGBoost's own near-zero R2 there is an independent reason
  to trust it less on that cohort specifically.
\item RQ1: the DNN/PINN comparison is cohort-conditional (DNN wins on 4survey, ties on
  6survey), alongside a real difference in the two cohorts' fitted growth curves.
\end{enumerate}

\section*{3. Ranked figure plan}

Five figures across three research questions, at roughly one per question with two
questions given a second panel where the underlying finding genuinely needs two
comparisons shown together. Every figure below already exists as a fully specified
entry in \texttt{d\_draft\_plan\_16th\_855pm\_plotting.md}; this list only selects and
re-ranks a subset against the "one to two, three at most, per question" budget.

\textbf{Essential (5 total)}

\begin{itemize}
\item \textbf{RQ1 -- paired fold comparison, XGBoost vs.\ DNN.} Per-fold R2 for both
  models on both cohorts, one line per fold connecting the two models' scores
  (paired-points design, not a bar chart), coloured by cohort. Establishes both the
  size of the XGBoost advantage \emph{and} whether it holds fold-by-fold in one image,
  which the results table alone cannot show (the table only gives the pooled mean and
  SD, not the paired direction per fold). Adds: the 4/5-fold and 5/5-fold win counts,
  and the one reversed fold (4survey, fold 0), directly.
\item \textbf{RQ2a -- where conditioning helps and hurts, by quartile.} Mean
  reduction per CR-baseline-error quartile (Q1--Q4), paired points with error bars
  (fold SD), one series per model (DNN, XGBoost), faceted by cohort. Shows the
  quartile-level trade-off the aggregate table hides: Q1 degradation and Q4 correction
  in the same figure, plus whether XGBoost and DNN agree at the quartile level, not
  only in aggregate.
\item \textbf{RQ2b -- attribution and its leftover spatial residual (composite).}
  Panel A: coefficient forest plot (NLME and EN side by side, one row per variable,
  point = mean, whisker = fold SD, faceted by set). Panel B: Moran's I by set for both
  EN and XGBoost, flat across sets. Panel C: a map of each model's leftover residual
  on Set4/4survey. Establishes the attribution ranking, the flat-clustering finding,
  and where the unexplained variance actually sits, together, since Panel C is what
  motivates RQ3 directly.
\item \textbf{RQ3 -- GNNWR's edge and its spatial meaning (composite).} Panel A:
  pooled R2 by set and model (EN/XGBoost/GNNWR) with bootstrap CIs, 4survey. Panel B:
  residual Moran's I by set/model/cohort, showing the 4survey reduction and 6survey
  increase side by side. Panel C: GNNWR's own per-plot local coefficient map for the
  variable ranked first. Establishes the accuracy edge, its CI-overlap caveat, its
  independent spatial corroboration, and what a genuinely spatially-varying
  coefficient looks like, in one figure.
\item \textbf{RQ3 -- where deviations occur and why (composite).} Panel A: a map of
  \texttt{local\_y\_max\_difference} across Aberfoyle. Panel B: compartment archetype
  map (the six classes from the archetype check). Panel C: representative trajectory
  plots, one panel per archetype, observed height vs.\ age against the fitted local
  curve and the yield-class benchmark. Establishes that the deviation is spatially
  structured, what kind of deviation each archetype represents, and what a
  representative trajectory actually looks like, which no table can show.
\end{itemize}

\textbf{Useful if space allows}

\begin{itemize}
\item RQ1 -- split-difficulty gradient (plot-level to spatial-block to temporal),
  dumbbell plot, one row per model. Reinforces the leakage-resistance argument
  visually; not essential since the plot-level/temporal numbers are already stated in
  prose with clear direction.
\item RQ3 -- CanopyCover rank agreement across EN/XGBoost/GNNWR as a rank
  slopegraph, 4survey vs.\ 6survey. Makes the 6survey model-split finding (headline
  \#9) visually obvious; useful but the finding is already fully stated in one
  sentence with exact ranks, so a reader does not strictly need the image.
\end{itemize}

\textbf{Appendix only}

\begin{itemize}
\item RQ1 -- two fitted CR curves (4survey vs.\ 6survey) with elevation
  distributions. Background/context for the cohort-difference explanation, not
  evidence for a specific claim.
\item RQ2b -- VIF diagnostic (lollipop plot). Supports the "multicollinearity ruled
  out" aside; not a headline finding.
\item RQ3 -- compartment archetypes' environmental means (bar chart, appendix
  version of the boundary-proximity and wind-exposure numbers already in prose).
\end{itemize}

\textbf{Omit}

\begin{itemize}
\item Any grouped-bar reproduction of the RQ1 results table (Figure R1 originally
  proposed as a plain grouped bar): the table already states these seven numbers per
  cohort clearly; a bar chart of the same six columns adds nothing the paired-fold
  figure above does not already show better.
\item RQ2b Set1-vs-Set4 ablation bar chart: the same two numbers (R2 0.220 vs.\ 0.358,
  Set1 vs.\ Set4) are clearer stated once in prose than as a two-bar chart.
\end{itemize}

\section*{5. Unsupported statements, missing evidence, and points needing checking}

\begin{enumerate}
\item \textbf{RQ1:} the specific forestry mechanism linking lower elevation to
  6survey's different curve shape (faster early growth, lower ceiling) is not
  established, only the compositional association. \chk{a domain explanation or a
  citation connecting elevation to this specific growth-curve shape difference}.
\item \textbf{RQ1:} "gradient-boosted trees beat neural networks on small tabular
  data" is a literature-level explanation (Grinsztajn et al.\ 2022), not something
  tested directly in this project. State it as a plausible general account, not a
  project finding.
\item \textbf{RQ2b:} no seed-sweep exists for RQ2b at all; every RQ2b number (the
  attribution table, coefficients, Moran's I) reflects split-seed 42 only. This is a
  standing limitation that should be stated plainly wherever RQ2b's numbers are
  presented as stable, not only in a footnote.
\item \textbf{RQ2b:} why NLME and Elastic Net disagree specifically on \texttt{topex}
  (stable in one, unstable in the other) is not tested; multicollinearity was checked
  and ruled out as the explanation, but no alternative explanation has been confirmed.
\item \textbf{RQ2b:} the static-climatology time-mismatch explanation for the unstable
  climate variables (\texttt{chelsa\_gdd5\_degc}, \texttt{tas\_mean}) is a disclosed
  property of the input data, not a tested cause; \texttt{soilgrids\_ph}'s own
  instability and \texttt{topex}'s EN-specific instability are not explained by this
  or any other hypothesis so far.
\item \textbf{RQ2b:} why canopy cover carries as much unique signal as the ablation
  shows (R2 drops by roughly a third when it is removed, more than its SHAP-rank
  margin alone would suggest) is not mechanistically explained beyond "it is the most
  direct structural measure available."
\item \textbf{RQ3:} the mechanism proposed for GNNWR's 6survey unreliability (too few
  reference compartments to estimate a stable local coefficient) is plausible and
  consistent with the reference-count numbers, but no ablation varying reference
  density was run to confirm it directly.
\item \textbf{RQ3:} whether letting the growth-rate parameter float (instead of fixing
  it at the yield-class value) would reduce the outlier-plot artifact is untested;
  RQ2a's own PINN\_k identifiability finding is the reason to expect it would not help
  cleanly, not direct evidence either way for RQ3's specific fitting procedure.
\item \textbf{RQ3:} the sub-compartment spatial clustering of the 266 flagged plots
  (item 7 in the source draft) currently has no confirmed-buildable figure; it would
  need real compartment-boundary polygon data not yet confirmed available.
\item \textbf{General, all of RQ2b and RQ3:} XGBoost's hyperparameter configuration is
  shared across RQ2b and RQ3 rather than tuned per target (only RQ1's XGBoost was
  individually tuned this pass). Any claim that leans on XGBoost's own attribution
  being as reliable as EN's or NLME's should carry this caveat, and RQ3's own 6survey
  finding (headline \#9 above) already uses this exact limitation as part of its own
  explanation.
\end{enumerate}

% ============================================================================
% END WORKING NOTES
% ============================================================================

\chapter{Results and Discussion}
\label{ch:results}

\section{RQ1: Top-height prediction}
\label{sec:results-rq1}

A tuned XGBoost model gives the most accurate top-height predictions of every model
tested, on both cohorts, under the spatial-block evaluation used throughout this
chapter (Table~\ref{tab:results-rq1}). Chapman-Richards guidance does not recover this
gap at any tested loss weight, though it buys back a specific, narrowly-scoped
property that only the guided architecture can produce in the first place.

\begin{table}[!htbp]
\centering
\scriptsize
\renewcommand{\arraystretch}{1.2}
\begin{tabular}{lrrrrrr}
\toprule
Model & 4survey RMSE & 4survey MAE & 4survey R2 & 6survey RMSE & 6survey MAE & 6survey R2 \\
\midrule
Linear & 5.168$\pm$0.361 & 4.056$\pm$0.343 & 0.580$\pm$0.026 & 4.877$\pm$0.643 & 3.619$\pm$0.426 & 0.509$\pm$0.067 \\
RF & 4.794$\pm$0.187 & 3.611$\pm$0.142 & 0.638$\pm$0.011 & 4.062$\pm$0.249 & 3.010$\pm$0.185 & 0.656$\pm$0.046 \\
XGBoost (raw defaults) & 4.797$\pm$0.175 & 3.621$\pm$0.136 & 0.637$\pm$0.015 & 4.106$\pm$0.358 & 3.041$\pm$0.301 & 0.651$\pm$0.016 \\
\textbf{XGBoost (tuned)} & \textbf{4.548$\pm$0.137} & \textbf{3.431$\pm$0.123} & \textbf{0.674$\pm$0.009} & \textbf{3.656$\pm$0.232} & \textbf{2.669$\pm$0.183} & \textbf{0.722$\pm$0.031} \\
DNN & 4.681$\pm$0.230 & 3.533$\pm$0.201 & 0.655$\pm$0.016 & 3.903$\pm$0.316 & 2.867$\pm$0.275 & 0.684$\pm$0.031 \\
PINN & 5.206$\pm$0.399 & 4.063$\pm$0.407 & 0.573$\pm$0.036 & 3.886$\pm$0.290 & 2.819$\pm$0.196 & 0.686$\pm$0.038 \\
PINN\_k & 5.197$\pm$0.402 & 4.046$\pm$0.421 & 0.575$\pm$0.037 & 3.895$\pm$0.282 & 2.816$\pm$0.196 & 0.684$\pm$0.038 \\
\bottomrule
\end{tabular}
\caption{Model comparison for top-height prediction, Set3, spatial\_block\_kfold,
per-fold mean $\pm$ SD across five folds. 4survey n=232{,}292 row-years (58{,}073
plots); 6survey n=82{,}614 row-years (13{,}769 plots).}
\label{tab:results-rq1}
\end{table}

XGBoost's R2 advantage over the strongest neural competitor, the DNN, is 0.019 on
4survey and 0.038 on 6survey. In RMSE terms this is a 2.8\% reduction on 4survey and a
6.3\% reduction on 6survey; against PINN, a 12.6\% RMSE reduction on 4survey and 5.9\%
on 6survey. The gap is not a single-run artefact: matching the two models fold by
fold, XGBoost wins 5 of 5 folds on 6survey, though fold 4's margin is a bare 0.0005 R2,
and 4 of 5 folds on 4survey, with DNN ahead on the remaining fold by 0.0068. This is
reported as an honest exception rather than smoothed into the aggregate, since the
paired comparison is what the mean and SD in Table~\ref{tab:results-rq1} cannot show
on their own. A 27-configuration hyperparameter search run specifically on this
target (270 fits across both cohorts, selected on validation R2, never test) placed
the winning configuration within 0.001--0.004 R2 of a naively borrowed configuration,
well inside the fold-to-fold SD, so the advantage is not an artefact of a lucky
hyperparameter choice.

Why trees outperform the neural architectures here is not established by this
project's own evidence. Gradient-boosted trees are documented in the wider machine
learning literature as outperforming neural networks on small-to-medium tabular data
with relatively few, mixed-type features \citep{grinsztajn2022}, which is consistent
with, though not confirmed by, the comparison run here.

Within the neural family, the better model depends on the cohort. DNN beats PINN and
PINN\_k on 4survey with non-overlapping 95\% confidence intervals (DNN
[0.624, 0.690] vs.\ PINN [0.545, 0.607]); on 6survey the two intervals overlap fully
(DNN [0.658, 0.740] vs.\ PINN [0.671, 0.736]), so this reads as a tie rather than a
loss, and the tie is stable across a five-seed reseed. Two compositional differences
accompany this: 6survey's own fitted Chapman-Richards curve has a 7.5\,m lower
asymptote (44.46\,m vs.\ 51.96\,m), a growth rate more than double 4survey's
(k=0.0233 vs.\ 0.0103), and a different shape ($p$=1.193 vs.\ 0.865, not simply a
rescaled version of the same curve); 6survey also samples a narrower, lower-elevation
slice of Aberfoyle (mean 102\,m vs.\ 178\,m, range capped at 351\,m vs.\ 561\,m). The
same asymmetry shows up when forecasting forward in time: PINN loses 50\% of its
single-split R2 moving from spatial-block to temporal evaluation on 4survey but 78\%
on 6survey, and DNN loses 14\% versus 56\%. This is consistent with, though not
separately proven beyond, a smaller and more homogeneous training population giving
every model less to generalise from, on both the spatial and temporal axes at once.

The choice of evaluation split changes the picture entirely for one model and barely
touches the other two. Under an easy plot-level split, where a held-out plot's nearest
neighbour is typically metres away in the training set, DNN's R2 inflates by 0.197 on
4survey and 0.122 on 6survey; PINN and PINN\_k inflate by 0.006 or less on either
cohort. The reason is architectural: DNN's environmental input feeds directly into the
same 128-unit trunk used for everything else, letting it fit an effectively
per-location signature from near-duplicate nearby plots, while PINN's terrain and wind
features only ever pass through a 16-unit sub-network producing a single additive
adjustment to the curve's asymptote, further constrained by the physics and trajectory
losses. This is a standard bias-variance argument, not a speculative one, and it is
the direct justification for using the spatial-block split as the primary evaluation
design throughout the rest of this chapter rather than the easier plot-level split.

The physics and trajectory losses carry a measurable accuracy cost. R2 falls
monotonically as their weight increases from 0 to 1 to 2, for every model and every
cohort. On 4survey the unconstrained model (weight 0) beats the default weight of 1 by
0.052 R2 (PINN) and 0.050 R2 (PINN\_k), clearly outside either configuration's own
confidence interval; on 6survey the same direction holds but the gap (0.010 for both
variants) sits inside 6survey's wider interval, so it is not a statistically clean
effect on that cohort alone. Removing the constraint does buy something back for
PINN\_k specifically: with the constraint in place, its learned asymptote and growth
rate are correlated at $-0.449$ on 4survey; without it, at $-0.935$, meaning the two
parameters trade off against each other almost perfectly rather than varying
independently. This is a standard growth-curve identifiability problem: without
survey data near a plot's true asymptote, a taller-but-slower and a shorter-but-faster
trajectory fit the observed heights almost equally well, and the physics constraint
resolves that ambiguity by anchoring the model to the curve's analytical derivative.
The constraint therefore trades a real amount of predictive accuracy for a real amount
of parameter identifiability, but only for PINN\_k's own two-parameter design; XGBoost
has no comparable parameters to lose or keep.

\textbf{For RQ1}, the evidence supports XGBoost as the most accurate top-height model
under spatial-block evaluation on both cohorts tested here, and supports the
spatial-block split itself as necessary rather than a formality, given how far the
plot-level split's numbers move for the one model architecturally able to exploit it.
It does not establish why tree ensembles win over neural networks on this dataset, and
the case for retaining Chapman-Richards guidance rests on the narrower claim that it
preserves a specific interpretability property PINN\_k would otherwise lose, not on
raw predictive accuracy.

\fig{RQ1 -- paired per-fold R2, XGBoost vs.\ DNN, both cohorts. See working notes \S3.}

% ----------------------------------------------------------------------------

\section{RQ2: Attribution to the shared growth curve}
\label{sec:results-rq2}

\subsection{Does environmental conditioning correct the shared curve's departures?}
\label{sec:results-rq2a}

A model given environmental information does not improve prediction uniformly
relative to the shared curve: it corrects the plots the curve already predicts worst,
and is noticeably less accurate on the plots it already predicts well
(Table~\ref{tab:results-rq2a}). Only the first half of this turns out to be an effect
of the environmental information itself.

\begin{table}[!htbp]
\centering
\scriptsize
\renewcommand{\arraystretch}{1.2}
\begin{tabular}{llrrrr}
\toprule
Model & Cohort & CR baseline error (m) & Mean reduction (m) & \% of baseline error & \% rows improved \\
\midrule
DNN, env-conditioned & 4survey & 5.028 & 1.501$\pm$0.023 & 29.9\% & 64.8$\pm$4.1 \\
DNN, env-conditioned & 6survey & 3.047 & 0.174$\pm$0.036 & 5.7\% & 52.4$\pm$4.8 \\
XGBoost, env-conditioned & 4survey & 5.028 & 1.597$\pm$0.330 & 31.8\% & 67.2$\pm$3.8 \\
XGBoost, env-conditioned & 6survey & 3.047 & 0.361$\pm$0.183 & 11.8\% & 56.0$\pm$4.6 \\
XGBoost, no-env control & 4survey & 5.030 & 1.222$\pm$0.232 & 24.3\% & 64.5$\pm$3.2 \\
XGBoost, no-env control & 6survey & 3.047 & 0.240$\pm$0.144 & 7.9\% & 54.8$\pm$3.0 \\
\bottomrule
\end{tabular}
\caption{Residual reduction against the fold-matched Chapman-Richards baseline
(Equation~\ref{eq:residual-reduction}), Set3, pooled across five spatial-block folds.
DNN row is a five-seed reseed; XGBoost rows are a single fixed configuration.}
\label{tab:results-rq2a}
\end{table}

The aggregate reduction (29.9\% of baseline error on 4survey, 5.7\% on 6survey for
DNN) hides a sharp split by how well the curve already fit a given plot. Splitting
test rows into quartiles by the curve's own baseline error, from Q1 (best-fitting
quartile) to Q4 (worst), shows DNN's error in Q1 is 2.6 times larger than the curve's
own Q1 error on 4survey and 3.2 times larger on 6survey; only 21--25\% of Q1 rows
improve at all. The correction is concentrated in Q4, where the absolute gain reaches
5.08\,m on 4survey and 1.81\,m on 6survey. This pattern replicates across every one of
18 model/set/cohort combinations tested and is stable across a five-seed reseed
(4survey mean reduction 1.501$\pm$0.023\,m; Q4 alone 5.054$\pm$0.117\,m, an SD under
2.5\% of the mean).

A plain XGBoost model, given the same environmental features and no physics guidance
at all, reproduces this pattern closely: its mean reduction (1.597\,m on 4survey,
0.361\,m on 6survey) is comparable to or larger than DNN's, its own Q1 relative-harm
ratio (2.1--2.4x) sits in the same range as DNN's, and its fold-to-fold SD is if
anything slightly smaller (69--85\% of DNN's own SD on both cohorts). Comparing the
environment-conditioned XGBoost against a no-environment control isolates what drives
the gain: the environmental features themselves add 0.375\,m on 4survey and 0.121\,m
on 6survey over the no-environment arm. Since a model with no physics loss and no
Chapman-Richards embedding shows the same quartile-concentrated pattern as the guided
neural model, the effect looks like a property of the environmental data correlating
with the curve's residual in this quartile-dependent way, rather than a property of
either model's architecture.

Seed-to-seed stability alone could not settle why the pattern splits this way: a
five-seed reseed showed Q4's effect is the most reproducible quartile across
independent seeds (coefficient of variation 2.3--4.5\%), consistent with real signal,
but Q1 was not the least reproducible quartile in either cohort, Q2 was, most likely
because Q2's near-zero mean inflates its own coefficient of variation mechanically
rather than because Q2 is genuinely noisier. Reproducibility across seeds cannot
distinguish a real signal being fit every time from the same noise pattern in the
training data being fit the same way every time. A permutation control resolves this
directly: refitting the same XGBoost model on Set3 with its environmental columns
shuffled across rows before training, so the model sees real heights but a fake,
scrambled environment. Q1's degradation is nearly unchanged under this fake
environment (4survey: $-1.30$\,m permuted vs.\ $-1.21$\,m real; 6survey: $-0.57$\,m
permuted vs.\ $-0.72$\,m real), so the degradation does not require real
environmental signal at all; it is a property of adding a flexible model to rows the
baseline already fit well. Q4's correction shrinks by 18\% (4survey) to 27\%
(6survey) once the environment is permuted away, so that half of the pattern is a
genuine, if partial, effect of real environmental data.

\textbf{For RQ2's first part}, the evidence supports the hypothesis that real
environmental information corrects the plots the shared curve predicts worst, at
least partially. It does not support treating the accuracy loss on well-fitting plots
as a cost of environmental conditioning specifically: that loss appears at a similar
size even when the model is given scrambled, meaningless environmental values, so it
reflects a cost of model flexibility in general, not of using real environmental
data.

\subsection{What explains a plot's persistent departure from the shared curve?}
\label{sec:results-rq2b}

Restricting to the 4survey cohort and modelling each plot's mean departure across its
full survey history directly (Table~\ref{tab:results-rq2b}), canopy cover and thinning
history dominate the attribution in every set and every method tested.

\begin{table}[!htbp]
\centering
\scriptsize
\renewcommand{\arraystretch}{1.2}
\begin{tabular}{lrlll}
\toprule
Set & n rows & EN R2 [95\% CI] & XGB R2 [95\% CI] & NLME spatial variance explained \\
\midrule
Set1 (baseline only) & 71{,}766 & 0.220 [0.180, 0.252] & 0.222 [0.183, 0.252] & 0.016$\pm$0.078 \\
Set2 & 71{,}766 & 0.359 [0.310, 0.396] & 0.416 [0.355, 0.461] & 0.204$\pm$0.066 \\
Set3 & 71{,}727 & 0.325 [0.272, 0.365] & 0.375 [0.321, 0.414] & 0.049$\pm$0.048 \\
Set4 & 71{,}330 & 0.358 [0.305, 0.398] & 0.395 [0.338, 0.438] & 0.175$\pm$0.064 \\
\bottomrule
\end{tabular}
\caption{Attribution of the persistent mean curve residual (Equation~\ref{eq:mean-cr-residual}),
4survey only, spatial-block split.}
\label{tab:results-rq2b}
\end{table}

Canopy cover carries the largest or near-largest Elastic Net coefficient in every set
(Set4: +1.943$\pm$0.034) and ranks first by mean $|$SHAP$|$ in every XGBoost set
(1.375--1.442). Because canopy cover is drawn from the same LiDAR pipeline as the
height targets themselves, its correlation with all three of this project's own
targets was checked directly rather than assumed benign: Pearson correlations of
0.35--0.47, moderate rather than the target restated in different units. Removing
canopy cover from Set4 and refitting drops R2 by roughly a third for both models (EN:
0.350 to 0.231, XGBoost: 0.388 to 0.249), a larger effect than its SHAP margin alone
(1.2--1.7x over the next variable) would suggest, meaning it carries signal nothing
else in the set backfills; the variable most likely to absorb the gap,
\texttt{dist\_to\_road}, does not (its own SHAP value moves slightly the wrong way,
1.003 to 0.869), which rules out simple correlation-driven credit transfer as the
mechanism without explaining why canopy cover carries this much unique signal in the
first place. The defensible reading is to treat canopy cover and thinning as baseline
stand-structure controls, expected to dominate by construction, and to look at the
true environmental variables' contribution beyond that baseline as the actual
attribution question.

Among the abiotic environmental variables, only \texttt{slope\_degrees} is reliably
stable across model families: NLME gives +0.894/+0.861 and Elastic Net +1.230/+0.993
on Set3/Set4, the same sign and comparable magnitude in both. \texttt{topex} is
strong and consistent within NLME (+1.316/+1.323) but weak and sign-flipping within
Elastic Net ($-0.175$ Set3, $+0.063$ Set4, both with an SD larger than or comparable
to the coefficient itself), a genuine disagreement between the two methods rather than
a shared, confirmed signal. Multicollinearity was checked directly as a possible
explanation and ruled out: the unstable variables (\texttt{tas\_mean}, VIF 3.44;
\texttt{soilgrids\_ph}, VIF 2.12) do not carry higher VIF than the stable
\texttt{chelsa\_bio12\_precip\_mm} (VIF 3.91), and \texttt{topex} itself (VIF 2.37) is
lower than several stable variables. A static 1981--2010 climatology applied
identically regardless of a plot's actual 2008--2023 survey year is a disclosed
property of the climate inputs and a plausible contributor to the climate variables'
instability, though this has not been isolated as the confirmed cause, and it does not
explain \texttt{soilgrids\_ph} (a genuinely static soil property) or \texttt{topex}'s
own EN-specific instability at all.

Spatial clustering in the model residuals barely moves as more environmental data is
added. Moran's I stays in the 0.65--0.74 range from Set1 through Set4 for both EN and
XGBoost (p=0.001 throughout); going from four baseline columns to nineteen
VIF-screened columns across every category nudges EN's own Moran's I from 0.738 to
only 0.693. Canopy cover and thinning explain a real share of row-level variance, but
most of the spatially-structured variance in the departure is left unexplained
regardless of how much environmental data is supplied. Every set's confidence
interval clears the baseline-only Set1 interval with no overlap, so environmental data
does add real value over stand structure alone; Set2 and Set4's own intervals overlap
substantially ([0.355, 0.461] vs.\ [0.338, 0.438] for XGBoost), so which of the two is
genuinely better is not a statistically clean claim, while Set3, terrain-heavy with no
soil or climate coverage, is more clearly the weakest set on point estimates and
interval position alike.

\textbf{For RQ2's second part}, structural and management variables explain most of
the attributable signal in a plot's persistent departure from the shared curve; among
true environmental variables, only slope shows a reliably consistent association
across independent model families, and a substantial share of the departure's spatial
pattern is not explained by any environmental set tested here, which is the direct
motivation for testing a spatially-varying model in RQ3.

\fig{RQ2a -- quartile reduction by model and cohort. RQ2b -- coefficient/Moran's
I/residual-map composite. See working notes \S3.}

% ----------------------------------------------------------------------------

\section{RQ3: Plot-specific curve deviation}
\label{sec:results-rq3}

RQ2 established that global attribution models leave substantial spatial clustering
unexplained (Moran's I 0.65--0.74) regardless of how much environmental data they are
given. RQ3 tests directly whether letting the environment-deviation relationship vary
spatially, rather than fitting one global relationship, captures what the global
models could not (Table~\ref{tab:results-rq3}).

\begin{table}[!htbp]
\centering
\scriptsize
\renewcommand{\arraystretch}{1.2}
\begin{tabular}{lllllll}
\toprule
Set & EN 4survey [CI] & EN 6survey & XGB 4survey [CI] & XGB 6survey & GNNWR 4survey [CI] & GNNWR 6survey \\
\midrule
Set2 & 0.286 [0.241, 0.325] & 0.057 & 0.266 [0.214, 0.308] & 0.031 & 0.320 [0.270, 0.369] & $-0.044$ / 0.014$\pm$0.053 \\
Set3 & 0.274 [0.223, 0.316] & 0.031 & 0.281 [0.235, 0.320] & 0.086 & 0.319 [0.263, 0.365] & 0.058 / 0.054$\pm$0.055 \\
Set4 & 0.240 [0.192, 0.286] & 0.056 & 0.250 [0.201, 0.295] & 0.015 & 0.294 [0.236, 0.347] & $-0.075$ / 0.002$\pm$0.068 \\
\bottomrule
\end{tabular}
\caption{Attribution of the plot-specific asymptote deviation
(Equation~\ref{eq:local-ymax-diff}), spatial-block split. 6survey GNNWR column shows
seed-42 point estimate / three-seed mean$\pm$SD.}
\label{tab:results-rq3}
\end{table}

GNNWR's point estimate beats both global models on every 4survey set, consistently in
direction and magnitude rather than as a one-set fluke, but the advantage is not
confirmed at the strictest bar: EN and XGBoost's own bootstrap confidence intervals
overlap GNNWR's in every set (Set4: GNNWR [0.236, 0.347] vs.\ EN [0.192, 0.286]).
An independent diagnostic corroborates the same direction as the R2 gap without
resolving the interval overlap: on 4survey, GNNWR reduces residual Moran's I against
the best global model in all three sets ($-0.012$ to $-0.020$), evidence that its
advantage reflects capturing real spatial structure rather than a better mean fit
alone. On 6survey the same diagnostic reverses: GNNWR's R2 sign-flips across seeds
(mean near zero, within the seed-to-seed SD) and its residual Moran's I \emph{increases}
in all three sets ($+0.031$ to $+0.052$). Geographically weighted methods need enough
nearby reference points to estimate a stable local coefficient, and 6survey has far
fewer compartments than 4survey (47 vs.\ 231); this is a plausible account for the
reversal, consistent with the reference-count difference, though no ablation varying
reference density was run to confirm it as the specific cause.

Which variable ranks first also depends on cohort, and on 6survey specifically it
splits by model rather than by attribution method. On 4survey, canopy cover ranks
first by three converging methods: the Elastic Net coefficient (2.3x the next-largest
variable on Set2 and Set4, 1.5x on Set3), XGBoost gain importance (consistently
2.6--2.7x across sets), and SHAP. On 6survey, EN's own coefficient still ranks canopy
cover first in all three sets, and GNNWR's own per-plot local coefficients
independently agree; XGBoost, by gain importance and SHAP alike, does not rank it
first in any 6survey set, beaten instead by \texttt{cpmt\_compactness\_ratio} or by
thinning-related variables depending on the set. XGBoost's own 6survey R2 for these
sets is close to zero or negative (Table~\ref{tab:results-rq3}), which is an
independent, non-circular reason to trust its 6survey attribution less: there is very
little genuine signal in that fit for gain importance or SHAP to attribute in the
first place. XGBoost was not retuned specifically for 6survey; this matches the
project's wider decision not to retune RQ2b or RQ3's shared XGBoost configuration at
all, and 6survey is not the primary cohort for RQ3's spatial questions regardless,
given GNNWR's own already-established unreliability there.

A separate line of evidence concerns which individual plots are hardest to fit. The
same roughly ten plots recur as the worst residual across nearly every feature set and
cohort combination, which rules out a poor choice of features as the explanation, and
none of them are disturbance-related: zero of ten trigger the project's own
clearfell, measurement-anomaly, or trajectory-instability flags. Half show a smooth,
internally consistent height trajectory whose large fitted deviation comes entirely
from that trajectory sitting far from its yield-class benchmark curve, consistent with
a probable yield-class lookup mismatch rather than any problem with the observed
heights; the other half show real, if sub-threshold, curve-fit instability across
their own survey years. Tracing five of the ten plots' biologically implausible fitted
asymptotes (up to 116\,m, against entirely normal raw observed heights of 38.9 to
46.8\,m) back to the fitting procedure itself confirms the mechanism directly: the
closed-form fit holds the growth-rate parameter fixed at the plot's yield-class value
and solves only for the asymptote, so a plot whose real trajectory is flatter or
slower than its assigned yield class assumes is forced into an inflated ceiling to
reconcile the mismatch. This is a property of how the target itself is constructed,
not a data-quality gap; no height- or asymptote-based filter exists in the cleaning
pipeline that would have caught it, since the raw heights are unremarkable. These same
plots sit roughly three times closer to compartment or block boundaries than a typical
plot (median 9.4\,m vs.\ 28.0\,m population-wide), a checkable, though not causally
confirmed, corroborating pattern consistent with boundary-adjacent plots being more
likely to have their own distinct trajectory than the single yield-class value applied
to the whole compartment assumes.

The pattern is not a handful of edge cases. A population-wide Tukey-fence check on the
fitted asymptote flags 0.47\% of 4survey (266 plots) and 5.32\% of 6survey (709
plots), clustered by compartment rather than scattered (one compartment alone has 12
of its 37 plots flagged). Dropping these plots from already-saved predictions barely
changes 4survey's pooled R2 ($-0.002$) but makes 6survey's R2 worse, not better
(EN $-0.027$, XGBoost $-0.037$, pushing XGBoost negative), so excluding them would
remove signal the models currently rely on rather than clean the data. Re-checking
the disturbance flags at this larger, statistically meaningful scale confirms the
earlier ten-plot finding holds: 0 of 266 clearfell-like, 0 of 266
measurement-inconsistent, only 2 of 266 ambiguous. A known artefact contributes a real
but minor share: 75.6\% of flagged plots have their single worst-fitting survey year
in 2008, well above the 47.8\% population rate, consistent with a documented
compartment-boundary redraw between the 2006 and 2008 surveys; dropping each plot's
2008 point and refitting moves 89\% of plots closer to their yield-class benchmark,
but by only around 6\% on average, a real but minor contribution rather than the main
cause. Flagged plots are, if anything, slightly more sheltered than typical, which
does not support windthrow as an explanation.

\textbf{For RQ3}, the evidence supports the hypothesis that the environment-deviation
relationship varies spatially, on the cohort with enough spatial reference density to
test it, though the accuracy edge alone does not clear the strictest statistical bar
and needs the independent Moran's I result to be read as a corroborated rather than a
confirmed finding. Separately, and independent of which attribution model is used, a
substantial share of the largest fitted deviations trace mechanically to the
target's own construction under a fixed growth-rate assumption rather than to a
confirmed ecological cause, which should temper how far any single extreme
\(\Delta y_{\max}^{\text{plot}}\) value is read as a real biological signal.

\fig{RQ3 -- GNNWR edge/spatial-structure/local-coefficient composite. RQ3 --
deviation map/archetype map/trajectory composite. See working notes \S3.}

\end{document}
